% GENERATED FILE. Edit canonical lessons, then run npm run generate:books.
% canonical-source-hash: fnv1a64:9906508e0947e42d
% canonical-lessons: RU-C438-ukazatel, RU-C438-tablichka, RU-C438-strelka, RU-C438-oriyentir, RU-C438-skhema

\chapter{Signs and landmarks}
\label{ch:ru-signs-and-landmarks}
\hlchaptermodality{438}

\begin{chapteropening}
\textbf{By the end of this chapter:} \emph{I can follow signs: a signpost, a sign on a door, an arrow, a landmark and a route map.}

Complete the last lesson of chapter 438: I can follow signs: a signpost, a sign on a door, an arrow, a landmark and a route map.
\end{chapteropening}

\section[\texorpdfstring{\ru{указатель} (ukazátel)}{указатель (ukazátel)}]{\ru{указатель} (ukazátel) — a signpost; an index}

\label{lesson:RU-C438-ukazatel}



\begin{quote}
Before the new one: say the Russian for a compass, then the Russian for a satnav.
\end{quote}

\subsection*{You'll want to know: \ru{указатель}}
\textbf{\ru{указатель}} (\emph{ukazátel}) — \textquotedblleft{}a signpost; an index\textquotedblright{}.

\begin{quote}
\ru{По} \ru{указателю} \ru{до} \ru{озера} \ru{три} \ru{километра}.
\end{quote}

\begin{quote}
\emph{Po ukazátelyu do ózera tri kilométra.}
\end{quote}

\begin{quote}
According to the signpost it is three kilometres to the lake.
\end{quote}

Masculine. \emph{\ru{Дорожный} \ru{указатель}} — a road sign.

\begin{cousinweb}
\emph{\ru{указатель}} is built on \emph{\ru{указать}} 'to show the way' with \emph{-\ru{тель}}, from \emph{\ru{казать}} 'to show', the root of \emph{\ru{сказать}} 'to say'. EtymDB lists \emph{\ru{указатель}} without a deeper source.
\end{cousinweb}

\begin{grammarlens}[title={What you've built}]
One more word for telling, explaining and describing at B1.
\end{grammarlens}

\subsection*{Your turn}
\begin{itemize}
\raggedright
  \item \emph{Say it:} \emph{ukazátel}
  \item \emph{Say it:} \emph{ukazátel}, once more
  \item \emph{Say it:} a sentence of your own with \emph{ukazátel}
  \item \emph{Recall:} say the Russian for a compass, then the Russian for a satnav, then say \emph{ukazátel} again
\end{itemize}

\begin{tcolorbox}[breakable,colback=teal!4,colframe=teal!35!black,title={Before you move on}]
What does \ru{указатель} mean? (\textquotedblleft{}A signpost; an index\textquotedblright{}.) Say it once more.
\end{tcolorbox}
\section[\texorpdfstring{\ru{табличка} (tablíchka)}{табличка (tablíchka)}]{\ru{табличка} (tablíchka) — a sign, a plaque, a nameplate}

\label{lesson:RU-C438-tablichka}



\begin{quote}
Before the new one: say the Russian for a satnav, then the Russian for a signpost.
\end{quote}

\subsection*{You'll want to know: \ru{табличка}}
\textbf{\ru{табличка}} (\emph{tablíchka}) — \textquotedblleft{}a sign, a plaque, a nameplate\textquotedblright{}.

\begin{quote}
\ru{На} \ru{двери} \ru{висела} \ru{табличка} «\ru{Закрыто}».
\end{quote}

\begin{quote}
\emph{Na dverí viséla tablíchka «Zakrýto».}
\end{quote}

\begin{quote}
A sign saying Closed hung on the door.
\end{quote}

Feminine; a diminutive of \emph{\ru{таблица}}.

\begin{cousinweb}
\emph{\ru{табличка}} is a diminutive of \emph{\ru{таблица}} (EtymDB lists it without a source), which is usually traced to Latin \emph{tabula} 'a board, a plank' (Lewis \& Short).
\end{cousinweb}

\begin{grammarlens}[title={What you've built}]
One more word for telling, explaining and describing at B1.
\end{grammarlens}

\subsection*{Your turn}
\begin{itemize}
\raggedright
  \item \emph{Say it:} \emph{tablíchka}
  \item \emph{Say it:} \emph{tablíchka}, once more
  \item \emph{Say it:} a sentence of your own with \emph{tablíchka}
  \item \emph{Recall:} say the Russian for a satnav, then the Russian for a signpost, then say \emph{tablíchka} again
\end{itemize}

\begin{tcolorbox}[breakable,colback=teal!4,colframe=teal!35!black,title={Before you move on}]
What does \ru{табличка} mean? (\textquotedblleft{}A sign, a plaque, a nameplate\textquotedblright{}.) Say it once more.
\end{tcolorbox}
\section[\texorpdfstring{\ru{стрелка} (strélka)}{стрелка (strélka)}]{\ru{стрелка} (strélka) — an arrow (on a sign); a hand (of a clock)}

\label{lesson:RU-C438-strelka}



\begin{quote}
Before the new one: say the Russian for a signpost, then the Russian for a sign.
\end{quote}

\subsection*{You'll want to know: \ru{стрелка}}
\textbf{\ru{стрелка}} (\emph{strélka}) — \textquotedblleft{}an arrow (on a sign); a hand (of a clock)\textquotedblright{}.

\begin{quote}
\ru{Стрелка} \ru{на} \ru{знаке} \ru{показывает} \ru{дорогу} \ru{налево}.
\end{quote}

\begin{quote}
\emph{Strélka na znáke pokázyvayet dorógu nalévo.}
\end{quote}

\begin{quote}
The arrow on the sign shows the way to the left.
\end{quote}

Feminine. \emph{\ru{Часовая} \ru{стрелка}} — the hour hand; \emph{\ru{по} \ru{часовой} \ru{стрелке}} — clockwise.

\begin{cousinweb}
\emph{\ru{стрелка}} is a diminutive of \emph{\ru{стрела}} 'an arrow', which is inherited from Proto-Slavic \emph{strěla} 'an arrow' (EtymDB). EtymDB lists \emph{\ru{стрелка}} itself without a source.
\end{cousinweb}

\begin{grammarlens}[title={What you've built}]
One more word for telling, explaining and describing at B1.
\end{grammarlens}

\subsection*{Your turn}
\begin{itemize}
\raggedright
  \item \emph{Say it:} \emph{strélka}
  \item \emph{Say it:} \emph{strélka}, once more
  \item \emph{Say it:} a sentence of your own with \emph{strélka}
  \item \emph{Recall:} say the Russian for a signpost, then the Russian for a sign, then say \emph{strélka} again
\end{itemize}

\begin{tcolorbox}[breakable,colback=teal!4,colframe=teal!35!black,title={Before you move on}]
What does \ru{стрелка} mean? (\textquotedblleft{}An arrow (on a sign); a hand (of a clock)\textquotedblright{}.) Say it once more.
\end{tcolorbox}
\section[\texorpdfstring{\ru{ориентир} (oriyentír)}{ориентир (oriyentír)}]{\ru{ориентир} (oriyentír) — a landmark (to steer by)}

\label{lesson:RU-C438-oriyentir}



\begin{quote}
Before the new one: say the Russian for a sign, then the Russian for an arrow.
\end{quote}

\subsection*{You'll want to know: \ru{ориентир}}
\textbf{\ru{ориентир}} (\emph{oriyentír}) — \textquotedblleft{}a landmark (to steer by)\textquotedblright{}.

\begin{quote}
\ru{Хороший} \ru{ориентир} — \ru{церковь} \ru{на} \ru{площади}.
\end{quote}

\begin{quote}
\emph{Khoróshiy oriyentír — tsérkov na plóshchadi.}
\end{quote}

\begin{quote}
A good landmark is the church on the square.
\end{quote}

Masculine. Also a guideline: \emph{\ru{ориентир} \ru{для} \ru{работы}}.

\begin{cousinweb}
\emph{\ru{ориентир}} is built on \emph{\ru{ориентировать}} 'to orient', usually traced to French \emph{orienter} (EtymDB), on Latin \emph{oriēns} 'the east' (Lewis \& Short).
\end{cousinweb}

\begin{grammarlens}[title={What you've built}]
One more word for telling, explaining and describing at B1.
\end{grammarlens}

\subsection*{Your turn}
\begin{itemize}
\raggedright
  \item \emph{Say it:} \emph{oriyentír}
  \item \emph{Say it:} \emph{oriyentír}, once more
  \item \emph{Say it:} a sentence of your own with \emph{oriyentír}
  \item \emph{Recall:} say the Russian for a sign, then the Russian for an arrow, then say \emph{oriyentír} again
\end{itemize}

\begin{tcolorbox}[breakable,colback=teal!4,colframe=teal!35!black,title={Before you move on}]
What does \ru{ориентир} mean? (\textquotedblleft{}A landmark (to steer by)\textquotedblright{}.) Say it once more.
\end{tcolorbox}
\section[\texorpdfstring{\ru{схема} (skhéma)}{схема (skhéma)}]{\ru{схема} (skhéma) — a route map, a diagram}

\label{lesson:RU-C438-skhema}



\begin{quote}
Before the new one: say the Russian for an arrow, then the Russian for a landmark.
\end{quote}

\subsection*{You'll want to know: \ru{схема}}
\textbf{\ru{схема}} (\emph{skhéma}) — \textquotedblleft{}a route map, a diagram\textquotedblright{}.

\begin{quote}
\ru{У} \ru{входа} \ru{в} \ru{метро} \ru{висит} \ru{схема} \ru{линий}.
\end{quote}

\begin{quote}
\emph{U vkhóda v metró visít skhéma líniy.}
\end{quote}

\begin{quote}
There is a map of the lines at the metro entrance.
\end{quote}

Feminine. \emph{\ru{Схема} \ru{метро}} — the metro map.

\begin{cousinweb}
\emph{\ru{схема}} is usually traced to Greek \emph{skhêma} 'a form, a shape' (EtymDB lists the Greek word), probably through Latin \emph{schēma} or German \emph{Schema}; EtymDB does not list the Russian word.
\end{cousinweb}

\begin{grammarlens}[title={What you've built}]
One more word for telling, explaining and describing at B1.
\end{grammarlens}

\subsection*{Your turn}
\begin{itemize}
\raggedright
  \item \emph{Say it:} \emph{skhéma}
  \item \emph{Say it:} \emph{skhéma}, once more
  \item \emph{Say it:} a sentence of your own with \emph{skhéma}
  \item \emph{Recall:} say the Russian for an arrow, then the Russian for a landmark, then say \emph{skhéma} again
\end{itemize}

\begin{tcolorbox}[breakable,colback=teal!4,colframe=teal!35!black,title={Before you move on}]
What does \ru{схема} mean? (\textquotedblleft{}A route map, a diagram\textquotedblright{}.) Say it once more.
\end{tcolorbox}
